\documentclass[10pt]{article}
\usepackage[T1]{fontenc}
\usepackage[utf8]{inputenc}
\usepackage{lmodern}
\usepackage{booktabs}
\usepackage{graphicx}
\usepackage{hyperref}
\usepackage[margin=0.75in]{geometry}
\hypersetup{colorlinks=true,urlcolor=blue,pdftitle={Chicago Community Population - Dataset Nutrition Label}}
\pagestyle{empty}
\newlength{\NFwidth}
\setlength{\NFwidth}{5.00in}
\newcommand{\NFtitle}{\multicolumn{1}{c}{\fontsize{24}{26}\selectfont\bfseries Dataset Facts}}
\newcommand{\NFRULE}{\midrule[3.3pt]}
\newcommand{\NFrule}{\midrule[0.65pt]}
\begin{document}
\begin{center}
  % Self-contained label included by main.tex.
% Update these facts when the group changes the source CSV or analysis.
\begingroup
\renewcommand{\arraystretch}{1.08}
\setlength{\tabcolsep}{5pt}
\newcommand{\nfrow}[2]{\textbf{#1}\hfill #2\\}
\newcommand{\nfsection}[1]{\NFRULE\multicolumn{1}{l}{\scriptsize\bfseries #1}\\\NFrule}
\begin{tabular}{@{\hspace{6pt}}p{\NFwidth}@{\hspace{6pt}}}
\toprule
\NFtitle\\
\NFrule
\multicolumn{1}{c}{\large\bfseries Chicago Community Population}\\
\NFrule
\nfrow{Dataset size}{77 records \textbar\ 30 fields}
\nfrow{Record unit}{One Chicago community area}
\nfsection{Motivation}
\nfrow{Publisher}{Chicago Health Atlas}
\nfrow{Platform partners}{CDPH, UIC PHAME Center, Metopio}
\nfrow{Primary original source}{U.S. Census Bureau / ACS}
\nfrow{Research aim}{Demographics and air-pollution comparisons}
\nfsection{Composition}
\nfrow{Geographic scope}{All 77 Chicago community areas}
\nfrow{Reference label}{ACS Year 2023 (CSV field)}
\nfrow{Time window}{2019–2023}
\nfrow{Household income}{5 household-count brackets}
\nfrow{Age and sex}{12 population-count columns}
\nfrow{Race and ethnicity}{Counts; categories can overlap}
\nfrow{Missing cells}{0 null or blank cells found}
\nfsection{Collection}
\nfrow{Original collection}{American Community Survey}
\nfrow{Relevant source tables}{B01001 (age/sex), B03002 (race)}
\nfrow{Export type}{Custom indicator-selection CSV}
\nfsection{Cleaning and Labeling}
\nfrow{Rows removed}{None documented for population EDA}
\nfrow{Numeric parsing}{Removed commas; cast income counts}
\nfrow{Derived features}{Household totals and income shares}
\nfrow{Data-quality note}{Some subtotals differ by 1--2 people}
\nfsection{Uses and Distribution}
\nfrow{Intended uses}{Community-level comparisons}
\nfrow{Format/access}{Public CSV via Chicago Health Atlas}
\nfrow{Reuse}{Download, processing, redistribution reported}
\nfsection{Maintenance and Evolution}
\nfrow{Platform}{Atlas receives ongoing updates}
\NFRULE
\textbf{Responsible use:} Group patterns are not individual attributes;\\
2023-labeled demographics vs. 2025--2026 pollution data;\\
correlation does not demonstrate causation.\\
\bottomrule
\end{tabular}
\endgroup

\end{center}
\vspace{0.35em}
\noindent\begin{minipage}{0.96\linewidth}
\footnotesize
\textbf{Provenance:} Group's custom 30-column Chicago Health Atlas CSV (`Chicago\_Community\_Population\_Data.csv`), validated against the group's EDA documentation. POP/PCT/INC pages describe related indicators but do not verify every field in the custom export.\\[0.3em]
\textbf{Sources:} \href{https://chicagohealthatlas.org/}{Chicago Health Atlas}; \href{https://chicagohealthatlas.org/indicators/POP}{Population (POP)}; \href{https://chicagohealthatlas.org/indicators/PCT}{Demographics (PCT)}; \href{https://chicagohealthatlas.org/indicators/INC}{Income (INC)}.\\[0.3em]
\textbf{Reporting caution:} The source CSV says \texttt{ACS Year = 2023}, but its precise ACS multi-year reporting period and the release/version used for the custom export have not been established. The Atlas's current standalone POP export is \emph{not} numerically identical to the team's custom CSV.
\end{minipage}
\end{document}
